\subsection{Setup}
    \label{subsec:setup}
\paragraph{Datasets.}
    \label{para:datasets}
We evaluate our proposed model on two public echocardiography video datasets, CAMUS~\cite{leclerc2019deep} and EchoNet-Dynamic~\cite{ouyang2019echonet}. 
\textbf{CAMUS} contains 500 patient cases, each with apical four-chamber and two-chamber view videos. For each video, it provides frame-by-frame annotations spanning from end-diastole ($ED$) to end-systole ($ES$). These annotations include labeled contours of the left ventricular endocardium, left ventricular epicardium, and the left atrium. 
\textbf{EchoNet-Dynamic} contains 10,030 apical four-chamber view videos. It provides annotations for each video only at two key frames: $ED$ and $ES$. 
We used the original data splits for CAMUS and EchoNet-Dynamic, resizing them to resolutions of $256 \times 256$ and $128 \times 128$, respectively. From each cardiac video, we uniformly sampled 10 frames.

\noindent
\textbf{Evaluation metrics.}
    \label{para:eval}
To ensure a fair and comprehensive evaluation, we assessed the model's performance using both standard segmentation metrics and key clinical indices. The geometric accuracy of the segmentation was evaluated with four common metrics: the mean Dice coefficient (mDice), mean Intersection over Union (mIoU), Hausdorff Distance (HD), and Average Surface Distance (ASD). 
To evaluate the model's clinical utility for left ventricular analysis, we estimated the Left Ventricular Ejection Fraction ($LV_{EF}$) from the segmentation masks using Simpson's rule~\cite{LANG20151}. We then assessed the quality of these estimations by calculating the Pearson correlation coefficient (corr), mean bias (bias), and standard deviation (std) between the predicted and ground truth $LV_{EF}$ values.

\noindent
\textbf{Implementation details.}
    \label{para:imp}
All training runs on a single RTX 3090 GPU. 
For CAMUS and EchoNet-Dynamic, we conducted training over 1500 iterations. 
The AdamW optimizer~\cite{loshchilov2019decoupledweightdecayregularization} was used with a learning rate of 1e-4, a batch size of 10 and a variety of data augmentation techniques were applied, including gamma augmentation, random scaling, random rotation, and random contrast adjustment, each with a probability of 0.5. We clip the global gradient norm to $\lambda$ = 3 and use stable data augmentation. 
We use a combined loss function of cross-entropy and soft dice loss with equal weighting following \citet{cheng2024puttingobjectvideoobject}.